% !TeX spellcheck = en_US
\documentclass[french]{yLectureNote}

\title{Électromagnétisme 2 PS}
\subtitle{Physique}
\author{Paulhenry Saux}
\date{\today}
\yLanguage{Français}
%lionels.calmels@cemes.fr
\professor{M.Brut}
\usepackage{graphicx}%----pour mettre des images
\usepackage[utf8]{inputenc}%---encodage
\usepackage{geometry}%---pour modifier les tailles et mettre a4paper
%\usepackage{awesomebox}%---pour les boites d'exercices, de pbq et de croquis ---d\'esactiv\'e pour les TP de PC
\usepackage{tikz}%---pour deiffner + d\'ependance de chemfig
% \usepackage{tabularx}%---pour dimensionner automatiquement les tableaux avec variable X
\usepackage{awesomebox}%---Pour les boites info, danger et autres
\usepackage{menukeys}%---Pour deiffner les touches de Calculatrice
\usepackage{fancyhdr}%---pour les en-t\^ete personnalis\'ees
\usepackage{blindtext}%---pour les liens
\usepackage{hyperref}%---pour les liens (\`a mettre en dernier)
\usepackage{caption}%---pour la francisation de la l\'egende table vers Tableau
\usepackage{pifont}
\usepackage{array}%---pour les tableaux
\usepackage{lipsum}
\usepackage{yFlatTable}
\usepackage{multicol}
\usepackage{tabularx}
\usepackage{booktabs}
\newcommand{\Lim}[1]{\lim\limits_{\substack{#1}}\:}
\renewcommand{\vec}{\overrightarrow}
\newcommand{\N}[0]{\mathbb{N}}
\newcommand{\dd}{\mathrm{d}}
\newcommand{\norm}[1]{||\vec{#1}||}
\newcommand{\fo}{\psi(\vec{r},t)}
\newcommand{\foe}{\psi(\vec{r},t)\*}
\newcommand{\HH}{\hat{H}}
\newcommand{\hb}{\hbar}
\newcommand{\lap}{\nabla^2}
\newcommand{\lapcc}{\frac{\partial^2 }{\partial x^2}+\frac{\partial^2 }{\partial y^2}+\frac{\partial^2 }{\partial z^2}}
\newcommand{\E}{\vec{E}(M)}
\newcommand{\B}{\vec{B}(M)}
\newcommand{\V}{V(M)}
\newcommand{\er}{\vec{e_{\rho}}}
\newcommand{\ep}{\vec{e_{\varphi}}}
\newcommand{\pip}{\pi^+}
\newcommand{\pim}{\pi^-}
\newcommand{\mv}{\mathcal{V}}
\DeclareMathOperator{\Div}{Div}
\newcommand{\rot}{\vec{\mathrm{rot}}}
\newcommand{\grad}{\vec{\mathrm{grad}}}
\setcounter{chapter}{3}
\begin{document}
\chapter{Puissance et énergie du champ Em}
\section{Puissance cédée par le champ  à la matière}
Un champ électrostatique interagit avec des particules chargées et leur fournit de l'énergie.
Une particule chargée est soumise à la force de Lorenz :
$$\vec{F} = q(\vec{E}+\vec{v}\wedge \vec{B})$$ et la puissance de cette force :
$$P = \vec{F}\cdot \vec{v} = q\vec{E}\cdot \vec{v}+q(\vec{v}\wedge \vec{B})\cdot \vec{v} = q\vec{E}\cdot \vec{v}$$
On considère un volume élémentaire $\dd \tau$ qui contient une charge $\dd q$ égale à $ \rho(r,t) \dd \tau$
La puissance cédée par le champ au volume $\dd \tau$ :
$$\dd P = \dd q \vec{E}\cdot \vec{v} = \rho \dd \tau \vec{E}\cdot \vec{v} = \vec{j}\cdot \vec{E}\dd \tau$$
On définit la puissance volumique $\frac{\dd P}{\dd \tau} = \vec{j}\cdot \vec{E}$
En intégrant sur le volume, la puissance cédée à l'ensemble des charges vaut $$P = \iiint_{V} \vec{j}\cdot \vec{E}\dd \tau$$
\begin{theorem}[Puissance cédée à l'ensemble des charges]
    $$P = \iiint_{V} \vec{j}\cdot \vec{E}\dd \tau$$
\end{theorem}
\section{Bilan d'énergie électromagnétique}
\subsection{Démonstration du théorème de Poynting}
On développe le terme $\vec{j}\cdot \vec{E}$ en utilisant l'équation de Maxwell-Ampère.
$$\vec{\nabla}\wedge \vec{B} = \mu_0(\vec{j}+\epsilon_0 \frac{\partial \vec{E}}{\partial t})$$
donc 
$$\vec{j} = \frac{1}{\mu_0}\vec{\nabla}\wedge \vec{B}-\epsilon_0 \frac{\partial \vec{E}}{\partial t}$$
et 
$$ \vec{j}\cdot \vec{E} = \frac{1}{\mu_0}(\vec{\nabla}\wedge \vec{B})\cdot \vec{E}-\epsilon_0 \frac{\partial \vec{E}}{\partial t}\cdot \vec{E}$$
On étudie le premier terme :
$$ \vec{\nabla}\cdot (\vec{A}\wedge \vec{C}) = (\vec{\nabla}\wedge \vec{A})\cdot \vec{C}-\vec{A}\cdot (\vec{\nabla}\wedge \vec{c})$$
On se sert de cette formule pour obtenir :
$$ \frac{1}{\mu_0}(\vec{\nabla}\wedge \vec{B}\cdot \vec{E}) = \vec{\nabla}\cdot (\frac{\vec{B}}{\vec{E}}) + \frac{\vec{B}}{\mu_0}\cdot (\vec{\nabla}\wedge \vec{E})$$
On étudie ce deuxième terme :
$$\frac{\vec{B}}{\mu_0}\cdot (\vec{\nabla}\wedge \vec{E}) = \frac{\vec{B}}{\mu_0}\cdot(-\frac{\partial \vec{B}}{\partial t}) = -\frac{1}{\mu_0}\frac{\partial \norm{B}^2}{\partial t}$$
On obtient finalement pour le premier terme :
$$\frac{1}{\mu_0}(\vec{\nabla}\wedge \vec{B})\cdot \vec{E} = -\vec{\nabla}\cdot \frac{\vec{E}\wedge \vec{B}}{\mu_0}-\frac{1}{2\mu_0}\frac{\partial \norm{B}^2}{\partial t}$$
On étudie maintenant le deuxième terme :
$$-\epsilon_0 \frac{\partial \vec{E}}{\partial t}\cdot \vec{E} = -\frac{\epsilon}{2}\frac{\partial \norm{E}^2}{\partial t}$$
On obtient finalement le théorème de Poyting :
\begin{theorem}[théorème de Poyting]
    $$-\vec{j}\cdot \vec{E} = \vec{\nabla}\cdot (\frac{\vec{E}\wedge \vec{B}}{\mu_0})+\frac{\partial}{\partial t}(\frac{1}{2\mu_0}\norm{B}^2 + \frac{\epsilon_0}{2}\norm{E}^2)$$
\end{theorem}
\subsection{Analyse des termes}
\begin{itemize}
    \item Le premier terme est la puissance cédée par les charges au champ
    \item $\frac{\vec{E}\wedge \vec{B}}{\mu_0} = \vec{\Pi}$ est le vecteur de Poyting (densité de courant d'énergie en $W/m^2$)
    \item Le dernier terme est noté $\vec{u_{em}}$. C'est $\vec{u_e}+ \vec{u_m}$ qui est une densit volumique d'énergie électromagnétique.
\end{itemize}
On a donc une nouvelle écriture, la forme générale : 
\begin{theorem}[Théorème de conservation de l'énergie en présence d'une source]
    $$-\vec{j}\cdot \vec{E} = \Div\vec{\Pi}+\frac{\partial \vec{u_{em}}}{\partial t}$$
\end{theorem}
On a alors deux cas : 
\begin{itemize}
    \item Si $-\vec{j}\cdot \vec{E}>0$ : de l'énergie est crée, elle vient des charges.
    \item Si $-\vec{j}\cdot \vec{E}\neq 0$, il n'y a pas conservation de l'énergie magnétique
\end{itemize}
\begin{theorem}[Forme intégrale]
    $$\iint_V -\vec{j}\cdot \vec{E}\dd \tau = \iint \vec{\Pi}\cdot \vec{\dd S}+ \iint_\tau \frac{\partial \vec{u_{em}}}{\partial t}\dd \tau$$
    Le trosième terme est la variation d'énergie EM dans le volume, le deuxième le flux d'énergie à travers S et le premier terme est la puissance totale gagnée sur le volume.
\end{theorem}
La variation de l'énergie EM $\vec{u_{em}}$ dans un volume est due 
\begin{itemize}
    \item au travail des forces EM dans ce volume $(-\vec{j}\cdot \vec{E})$
    \item à un échange d'énergie avec l'extérieur $(\iint_S \vec{\Pi}\cdot \vec{\dd S})$
\end{itemize}
\section{Applications}
\subsection{Cas du condensateur avec stokage de l'énergie électrique}
On considère 2 armatures circulaires l'une au-dessus de l'autre de surface S séparées d'une distance e et de charge $\pm Q$.
On a une tension U entre les 2 disques.

On peut montrer que 
$$ \vec{E} = -\frac{\sigma}{\epsilon_0}\vec{e_z}$$ avec $\sigma = \frac{Q}{S}$
On a $\vec{E} = -\frac{Q}{\epsilon_0 S}\vec{e_z}$
On obtient la relation entre U et Q : $\frac{U}{e} = \frac{Q}{\epsilon_0 S}$

Donc $$Q = \frac{S\epsilon_0}{e} U = CU$$
On calcule maitenant $u_e$ :
$$u_e = \frac{\epsilon_0}{2}E^2 = \frac{Q^2}{2\epsilon_0 S^2}$$
Donc $U_{em} = \iint_{\tau}\dd \tau = \frac{Q^2}{2\epsilon S}e = \frac12 CU^2$
\subsection{Solénoide = Stokage de l'énergie magnétique}
On étudie un solénoide de longueur a<<l, parcouru par un courant I.

À l'intérieur du solénoide, on a un champ
$$\vec{B} = \mu_0n I\vec{e_z}$$On a 
$u_{em} = u_m = \frac{\mu_0 n^2I^2}{2} $. L'énergie totale vaut alors$U_{em} = \iint u_m \dd \tau = \frac{\mu_0 n^2I^2}{2}\Pi a^2 l$

$$U_{em} = \frac12 L I^2 $$ avec $L = \mu_0 n^2 \Pi a^ 2l$

\subsection{Résistance d'un cylindre conducteur}
Il est parcouru par un courant I et de conductivité $\gamma$ et on se place en base polaire.

On a $I =\iint \vec{j}\cdot \vec{\dd S}$ d'où $\vec{j} = \frac{I}{\pi a^2}\vec{e_z}$

On en déduit que $\vec{E} = \frac{\vec{j}}{\gamma} = \frac{I}{\gamma \pi a^2}\vec{e_z}$.

On cherche maintenant $\vec{B}$ : On suppose que l'on a un fil de longueur $l>>a$, donc $\B$ identique au cas du fil infini.

Donc on utilise Ampère : Avec un cercle de rayon $\rho$ centré sur l'axe du cylindre et $\B = B_{\theta}\vec{e_{\theta}}$

On a $\int B\cdot \vec{\dd l} = \mu_0 \iint \vec{j}\cdot \vec{\dd S} \Rightarrow B=2\pi\rho = \mu_0 j\pi \rho^2 \Rightarrow B(\rho) = \frac{\mu_0 j\rho}{2} = \frac{\my_0 I\rho }{2\pi a^2}$

Donc $\B = \frac{\mu_0 I\rho}{2\pi a^2}\vec{e_\theta}$

Exprimons le vecteur de Pointing :
$$\vec{\Pi} = \frac{1}{\mu_0}\E\wedge \B = \frac{1}{\mu_0} \frac{I}{\gamma \pi a^2}\vec{e_z}\wedge \frac{\mu_0 I\rho}{2\pi a^2}\vec{e_{\theta}}$$

Donc $\vec{\Pi} = - \frac{I^2 \rho}{2\pi a^4 \gamma}\vec{e_\rho}$

On calcule la densité d'énergie : $u_{em} = \frac{\epsilon}{2} \norm{E}^2 + \frac{1}{2\mu_0}\norm{B}^2 = \frac{\epsilon}{2}\frac{I^2}{\gamma^2\pi^2 a^4}+\frac{1}{2\mu_0}\mu^2I^2\rho^2\frac{1}{4\pi^2 a^4}$

Donc $\frac{\partial u_{em}}{\partial t}=0$ si I constant au cours du temps.

On doit avoir $-\vec{j}\cdot \E = \Div \vec{\Pi}$

Donc $-\gamma E^2 = -\frac{I^2}{\gamma (\pi a^2)^2} = \frac{1}{\rho} \frac{\partial}{\partial \rho}(\rho \pi_{\rho})$.

$\iint \vec{\pi}\cdot \vec{\dd S} =  \vec{\pi} \cdot (\rho\dd \theta \dd z)\vec{e_\rho}$ Les extrémités ne contribuent pas car $\vec{\pi}\perp \vec{e_z}$

On se place bien en $\rho=a$ et on obtient $-\frac{I^2}{2\pi^2 a^3 \gamma}2\pi al = -\frac{l}{\pi a^2 \gamma} = -R I^2$ : Pertes par effet Joules

La puissance est négative donc les charges ne contribuent à renforcer le champ.
\chapter{Ondes EM dans le vide}
\section{Équation de propagation}
\subsection{Équation de propagation des champs \E et \B}
On se place dans le vide et loin des sources.

$\Div \vec{E} = 0, \Div \vec{B} = 0, \rot \vec{E} = -\frac{\partial \vec{B}}{\partial t},\rot B = \mu_0 \epsilon_0 \frac{\partial \E}{\partial t}$

La Démonstration à partir de ces équations est à savoir faire.

\subsubsection{Champ électrique}
$$\rot(\rot E) = \grad(\div E)-\lap \vec{E}$$

$$\rot(-\frac{\partial \vec{B}}{\partial t}) = -\lap E, -\frac{\partial}{\partial t}(\rot B) = -\lap \vec{E}$$

Donc $ \lap \vec{E} = -\frac{1}{c^2}\frac{\partial^2 \vec{E}}{\partial t^2} = \vec{0}$

\subsubsection{Champ magnétique}
$$rot rot(B) = grad (div B) -\lap \vec{B}$$
$$\rot(\epsilon_0\mu_0 \frac{\partial E}{\partial t}) = -\lap \vec{B}$$
$$\epsilon_0\mu_0 \frac{\partial}{\partial t}(-\frac{\partial \B}{\partial t}) = -\lap B$$
donc $$\lap B - \frac{1}{c^2} \frac{\partial^2}{\partial t^2} = \vec{0}$$

Remarques :
\begin{itemize}
    \item On a un couplage des champs à l'orgine d'un phénomène de propagation d'une perturbation
    \item Une onde EM est une onde vectorielle à  3 dimensions. En optique ondulatoire, on ne conserve que le caractère vibratiire (modèle scalaire)
    \item Les équations de d'Alembert sont linéaires, donc les soltuions forment un sev. On peut regarder les phénomènes d'interférence : permet d'étudier des olutions monochormatique.
    \item 
\end{itemize}
\subsection{Équations de propagation des potentiels}
$$rot B = rot (rot A) = grad(div A)-lap A$$
$$rot B =1/c^2 \frac{\partial E}{\partial t} = 1/c^2 \frac{c^2\partial}{\partial t}(-grad V-\frac{\partial vec A}{\partial t}) = -grad (1/c^2 \frac{\partil v}{\partial t})-1/c^2 \frac{\partial \vec{A}}{\partial t}$$

d'où $grad (div A)-lap A  = -grad (1/c^2 \frac{\partial V}{\partial t})-1/c^2 \frac{\partial A^2}{\partial t^2}$

donc $\lap \vec{A} -1/c^2 \frac{\partial^2 A}{\partial t^2} = 0$ par la jauge de Lorenz.

On fait de meme pour V : div E=0, donc div(-gra dV-partial A/partial t) = 0, donc $\Delta V - 1/c^2 \frac{\partial^2 v}{\partial^2 t} = 0$

Jauge de Lorenz : Cndition pour que A et V se propagent comme E et B.

\section{Ondes planes progressives monochormatiques}
L'onde plane est sune solution unidimensionelle de l'équation d'onde. La propagation du champ EM ne dépend que d'une direction de l'espace
\subsection{Définitions}
\begin{theorem}[Surface d'onde]
    Surface sur laquelle la phase est uniforme à un temps donné, la surface sur laquelle le champ EM est uniforme
\end{theorem}
\begin{theorem}[Onde plane]
    Les surfaces d'onde sont des plans // entre eux. On a u.r l'équation des plans d'onde avec u la direction de propagation et r
\end{theorem}
\begin{definition}[Onde progressve]
    Plans d'onde avancent à une vitesse c. L'onde dépend de t-(u.r) /c
\end{definition}
\begin{definition}[Monochromatique]
    dépendance temporelle sinosodidale de pulsation w constante, varie en $\cos(\omega(t-\frac{u.r}{c})+\phi_0)$.
\end{definition}
Avec $\vec{k} = k\vec{u}, k = \omega/c, 2\pi/\lambda$
Donc $\cos(\omega t - k.r +\phi_0)$

Double périodicité temporelle et spatiale : $T = 2\pi/\omega, \lambda = 2\pi/k$

\subsection{Notation complexe}
Pour une OPPM : $\vec{E} = E_{0x} (\cos(\omega t-\vec{k}\cdot \vec{r})+\phi_x)+\dots$

En complexe : $\vec{E} = \vec{E_0}e^{-i\omega t - k.r}$ avec $\vec{E_0} = E_{0x}e^{-i\phi_x}+\dots$

$\vec{E} = Re(E(r,t))$
\subsection{Dérivation des OPPM}
$$E = E_0 e^{-i(\omega t)-(k_x x +\dots)}$$

Dérivée partielle par rapport au temps : 
$$\frac{\partial E}{\partial t} = -i\omega E$$

Nous allons monter que $\nabla = ik$

$$div E = \frac{\partial E_X}{\partial x} + \dots = E_x e^{-i(\omega t-k.r)}e^{-i\phi_x}(ik_x)+\dots = ie^{-i(\omega t-k.r)}(E_{0x}e^{-\phi_x}k_x+\dots)$$

Finalement, $div E = i\vec{k}\cdot \vec{E}$.

$rot E = i\vec{k}\wedge \vec{E}$.

On fait de meme avec B, ce qui nous permet de réécrire les éqauations de Maxwell : $i\vec{k}\cdot \vec{E} = 0, i\vec{k}\cdot \vec{B} = 0, i\vec{k}\wedge \vec{E}i\omeaga B, i\vec{k}\wedge \vec{B} = -i\omega 1/c^2 E


\end{document}
